\section{Literature review}
\label{sec:review}

Objective measurement of eating behaviour matters for obesity, eating
disorders and diabetes research, and self-report is known to be inaccurate
and cannot capture in-meal behaviour such as eating rate or pauses
\citep{diou2022intake}. Wearable sensors that detect chewing, bites or hand
gestures have therefore become a central approach, ranging from
classical machine learning on hand-crafted features to end-to-end deep
learning, with the aims of improving accuracy, reducing false detections
and generalising to people outside the training set \citep{bell2020scoping}.
\citet{fontana2014aim} introduced the Automatic Ingestion Monitor (AIM),
combining jaw-motion, hand-gesture and accelerometer sensors, and reported
89.8\% food-intake detection accuracy over 24 hours of free living. Its
successor, AIM-2 \citep{doulah2021aim2}, moved the sensors onto eyeglasses:
a temporalis-muscle sensor and a 3-axis accelerometer at 128\,Hz. It extracted
38 hand-crafted features per signal, selected nine by minimum-redundancy
maximum-relevance, and classified 10\,s epochs with a support vector machine,
reaching F1 0.818 under leave-one-subject-out validation on 30 participants,
as summarised in Table~\ref{tab:literature}. These feature-engineered models
work, but their performance is bounded by the features chosen by hand, and
they are sensitive to noise and to variation between people.

\begin{xltabular}{\textwidth}{@{}>{\raggedright\arraybackslash}p{30mm}
                               >{\raggedright\arraybackslash}X
                               >{\raggedright\arraybackslash}X
                               >{\raggedright\arraybackslash}X@{}}
\caption{State-of-the-art methods for food-intake and chewing detection, and
the two ECG studies whose methods transfer to this problem.}
\label{tab:literature}\\
\toprule
\rowcolor{black!8}
\textbf{References} & \textbf{Techniques} & \textbf{Goals} & \textbf{Findings} \\
\midrule
\endfirsthead
\toprule
\rowcolor{black!8}
\textbf{References} & \textbf{Techniques} & \textbf{Goals} & \textbf{Findings} \\
\midrule
\endhead
\bottomrule
\endlastfoot
\citet{doulah2021aim2} & AIM-2; hand-crafted features, mRMR selection, SVM &
  Food-intake detection with passive food-image capture &
  F1 0.818 (LOSO, 30 participants); camera triggered only during eating \\
\citet{diou2022intake} & Review: CNN on in-ear audio windows for chewing,
  smartwatch IMU for bites &
  Intake monitoring in free living &
  CNN F1 0.908 against 0.748 for an SVM; self-supervised learning comparable
  to supervised \\
\citet{kyritsis2021endtoend} & CNN + LSTM, end to end on raw wrist IMU &
  Bite and meal detection in the wild &
  Bite F1 0.923 (LOSO) \\
\citet{ghosh2022comparative} & Five deep networks (MLP, time-CNN, FCN,
  ResNet, Inception) on AIM-2 optical and accelerometer time series &
  Compare deep architectures for food-intake detection &
  ResNet best, balanced accuracy 93.47\% \\
\citet{ghosh2024integrated} & Wavelet scalogram of accelerometer, fused with
  food-image recognition &
  Reduce false positives in free living &
  Fused F1 80.77\%; fusion cuts false positives \\
\citet{stankoski2021smartwatch} & Selection of hard negative examples for
  training &
  Imbalanced data with imperfect labels &
  Person-independent F1 0.82 \\
\citet{wang2024eating} & Temporal CNN + multi-head self-attention &
  Eating speed from wrist IMU, towards free living &
  61 participants, 513\,h; eating-speed MAPE 0.110 \\
\citet{vedovelli2026recognition} & Transformer encoder and other models on
  smartwatch sensors &
  Eating-episode recognition &
  Transformer: sensitivity 0.691, specificity 0.504 (LOSO) \\
\citet{hannun2019cardiologist} & 34-layer 1D CNN on raw single-lead ECG &
  End-to-end arrhythmia classification &
  F1 0.837, above the average cardiologist (0.780) \\
\citet{ikram2025transformer} & Transformer on MIT-BIH beats &
  Five-class arrhythmia classification &
  Accuracy 97.1\%; split not at patient level \\
\end{xltabular}

\subsection{From hand-crafted features to learned representations}

To address the limits of hand-crafted features, several groups trained
networks directly on the signal. In this literature, a raw input denotes the
absence of hand-engineered features rather than of preprocessing:
z-normalisation is the only preprocessing in the convolutional time-series
baselines of \citet{wang2017strong}, and the end-to-end bite detector reviewed
by \citet{diou2022intake} receives preprocessed inertial signals.
\citet{papapanagiotou2017chewing} detected chewing with a 1D convolutional
network applied to in-ear audio windows; across window lengths of 1 to 5\,s,
the 5\,s network reached F1 0.908, against 0.748 for a support vector machine
on hand-crafted features, although such models ``require significantly
larger volumes of labeled data'' \citep{diou2022intake}. Similarly, \citet{kyritsis2021endtoend} trained a CNN feature extractor and
an LSTM jointly on raw wrist inertial data, reaching a bite-detection F1 of
0.923 under leave-one-subject-out evaluation. On AIM-2 itself,
\citet{ghosh2022comparative} compared five deep architectures on the optical
and accelerometer time series and found the residual network best, with a
balanced accuracy of 93.47\%. The same pattern holds outside eating: the
fully convolutional network of \citet{wang2017strong}, three blocks of
convolution, batch normalisation and ReLU, is a strong baseline across
time-series benchmarks.

The closest analogue to this problem is electrocardiography.
\citet{hannun2019cardiologist} classified arrhythmia with a 34-layer 1D CNN
applied end to end to raw single-lead ECG at 200\,Hz, and exceeded the
average cardiologist (F1 0.837 against 0.780). ECG and the AIM-2 signals
share a sampling rate of the same order, a quasi-periodic structure (heart
beats at about 1\,Hz, chews at 0.94--2.17\,Hz \citep{po2011chewing}) and the
need to recognise short events inside longer recordings. This motivates the
first and most important design choice in this work: the model reads the
normalised time series, not a spectrum. None of the reviewed studies
compares the two representations on the same network, so that comparison is
made directly in this work. \citet{ghosh2024integrated} took the opposite
route, converting accelerometer signals to a wavelet scalogram image. A
magnitude spectrum says which rhythms are present in a window but discards
when they occur, which is the information that separates a bite from
background motion; a scalogram keeps a time axis, but at a resolution fixed
by the transform rather than learned from the data.

\subsection{Attention and Transformers}

Attention mechanisms let a network weight the parts of its input that carry
evidence. Squeeze-and-excitation \citep{hu2018squeeze} gates feature channels
with a small learned function of their global average, at negligible cost.
Attention-based pooling \citep{ilse2018attention} treats a window as a bag
of instances and learns which instances determine its label; for eating,
the instances are time steps and a single bite may be the deciding one.
Transformers extend this to attention between all time steps.
\citet{ikram2025transformer} applied a Transformer to five-class arrhythmia
classification on MIT-BIH, reporting 97.1\% accuracy. The reported 80/20
split is of preprocessed samples rather than of patients, so the figure
measures performance on new beats from known patients, not on new
patients. In eating detection, \citet{wang2024eating} combined a temporal
CNN with multi-head self-attention to measure eating speed from wrist IMU
data on 61 participants, while \citet{vedovelli2026recognition} found that a
Transformer encoder on smartwatch data under leave-one-subject-out gave the
highest sensitivity (0.691) but a specificity of only 0.504, which is the
false-positive problem in its clearest form.

Despite their promise, Transformers carry many more parameters than the
attention layers above, and whether they help at a given data size is an
empirical question rather than a given. This work therefore uses the two
light attention mechanisms by default and treats a Transformer over the
convolutional features as a separate experiment, on the same folds.

\subsection{Supervision, augmentation and negative data}

Small annotated datasets call for more supervision per example and for more
examples per recording. Multi-task learning \citep{caruana1997multitask}
improves generalisation by training a shared representation on related
tasks; here, counting chews is the related task, since the ground truth
already contains chew annotations. This auxiliary task is specific to the
present work. Augmentation of wearable sensor data
raised accuracy from 77.54\% to 86.88\% in Parkinson's disease monitoring
\citep{um2017augmentation}; overlapping windows are the simplest form of it.
Window length itself trades resolution for context, and has a measurable
effect on recognition \citep{banos2014window}. For chewing, which ``typically
is in the range from 1 to 3\,Hz'', a 5\,s window was chosen so that ``even for
the lowest frequency of 1\,Hz, a window of 5\,s contains at least 5 periods''
\citep{diou2022intake}; the 8\,s windows used here contain at least eight.

False positives are the persistent weakness of eating detectors: many
non-eating activities produce jaw and head motion. Walking in particular has
a rhythm similar to chewing, 1 to 3\,Hz, and causes chewing classifiers ``to
yield high false-positive detections'', which accelerometer data can be used
to remove \citep{diou2022intake}. The same work notes that a detector can be
tuned ``towards higher recall or precision, or a balanced mode, based on the
needs of each use-case''. \citet{stankoski2021smartwatch}
reduced them by deliberately selecting hard negative examples for training,
and \citet{ghosh2024integrated} by fusing the sensor classifier with food-image
recognition. Finally, evaluation practice varies widely across the field
\citep{bell2020scoping}. \citet{saeb2017need} showed that record-wise
cross-validation, where windows from one person appear in both training and
test sets, can massively overestimate accuracy, and that subject-wise
validation is the one that matches the use case. The leave-one-subject-out
protocol used by \citet{doulah2021aim2} and \citet{diou2022intake} is adopted
here for the same reason.

\subsection{Contribution of this work}

\begin{itemize}
  \item A compact convolutional model with channel attention, temporal
        attention pooling and an auxiliary chew-count head, reading the
        normalised AIM-2 optical and accelerometer signals directly from the
        ETS database: 61 annotated sessions from 43 participants, evaluated
        at the participant level.
  \item A direct comparison of time-series and spectral input on the same
        network, which none of the reviewed studies reports. The time
        series reaches F1 0.858 against 0.796 and reduces false alarms from
        81 to 32 per hour of non-eating.
  \item A controlled ablation of each design decision, one at a time on
        identical participant-level folds, separating the components that
        carry the performance from those that do not.
  \item Negative results reported with their mechanism: temporal median
        smoothing and context windows without a target marker both lower
        F1, because meals at 8\,s resolution are punctuated rather than
        contiguous.
  \item A measured check of sensor-annotation alignment on every session,
        replacing an assumed correction with a verified one.
\end{itemize}

Despite the progress in deep learning for intake monitoring, several
challenges persist: limited annotated data, variation between participants
in sensor fit and chewing style, false positives from non-eating jaw and head
motion, and inconsistent evaluation that makes published results hard to
compare. Addressing these requires models that generalise to people they
have not seen, evaluated in the way they would be used.
